% !TEX program = xelatex
\documentclass[12pt, a4paper]{article}
\usepackage[utf8]{inputenc}
\usepackage[T1]{fontenc}
\usepackage[spanish]{babel}
\usepackage{xeCJK}
\xeCJKsetup{AutoFakeBold=true}
\setCJKmainfont{Noto Serif SC}
\setCJKsansfont{Noto Serif SC}
\setCJKmonofont{Noto Serif SC}
\usepackage{geometry}
\geometry{margin=2.5cm}
\usepackage{titlesec}
\usepackage{enumitem}
\usepackage{booktabs}
\usepackage{array}
\usepackage{xcolor}
\usepackage{mdframed}
\usepackage{parskip}
\usepackage{fancyhdr}
\usepackage{multicol}

\definecolor{azul}{RGB}{30, 90, 160}
\definecolor{grisclaro}{RGB}{245, 245, 245}
\definecolor{azulclaro}{RGB}{220, 235, 255}
\definecolor{rojosuave}{RGB}{200, 50, 50}

\titleformat{\section}{\large\bfseries\color{azul}}{}{0em}{}[\titlerule]
\titleformat{\subsection}{\normalsize\bfseries\color{azul!80}}{}{0em}{}

\setlength{\headheight}{14.5pt}
\addtolength{\topmargin}{-2.5pt}
\pagestyle{fancy}
\fancyhf{}
\rhead{\textcolor{azul}{\textbf{Chino Mandarín — Unidad 3}}}
\lhead{\textcolor{gray}{Guía 1}}
\rfoot{\thepage}
\lfoot{\textcolor{gray}{\small Lecciones 7 y 8 · Fechas, deseos y compras}}

\newmdenv[backgroundcolor=azulclaro, linecolor=azul, linewidth=1pt, innerleftmargin=10pt, innerrightmargin=10pt, innertopmargin=8pt, innerbottommargin=8pt]{cuadroazul}
\newmdenv[backgroundcolor=grisclaro, linecolor=gray!40, linewidth=0.5pt, innerleftmargin=10pt, innerrightmargin=10pt, innertopmargin=8pt, innerbottommargin=8pt]{cuadrogris}

\begin{document}

\begin{center}
{\LARGE \textbf{\textcolor{azul}{Guía de Estudio 1}}}\\[4pt]
{\large Lecciones 7 y 8}\\[2pt]
{\normalsize Tiempo estimado: 3--4 horas}
\end{center}

\vspace{0.5em}
\hrule height 2pt
\vspace{1em}

%---------------------------------------------------------
\section{Contenido teórico — Lección 7: Fechas y tiempo}
%---------------------------------------------------------

\subsection{1.1 Estructura de fechas}

En chino la fecha se expresa de mayor a menor: \textbf{año → mes → día → día de la semana}.

\begin{cuadroazul}
\textbf{Pregunta por el día del mes:} \quad 今天几号？\quad (jīntiān jǐ hào) — ¿Qué día es hoy?\\[4pt]
\textbf{Respuesta:} \quad 今天九月一号。\quad (jīntiān jiǔ yuè yī hào) — Hoy es el 1 de septiembre.\\[6pt]
\textbf{Pregunta por el día de la semana:} \quad 今天星期几？\quad (jīntiān xīngqī jǐ) — ¿Qué día de la semana es hoy?\\[4pt]
\textbf{Respuesta:} \quad 星期三。\quad (xīngqī sān) — Miércoles.
\end{cuadroazul}

\vspace{0.5em}
\noindent\textbf{Vocabulario de tiempo (hoy, ayer, mañana):}

\begin{center}
\begin{tabular}{llll}
\toprule
\textbf{Carácter} & \textbf{Pinyin} & \textbf{Significado} & \textbf{Ejemplo} \\
\midrule
今天 & jīntiān & hoy & 今天是星期三 \\
昨天 & zuótiān & ayer & 昨天是星期二 \\
明天 & míngtiān & mañana & 明天是星期四 \\
\bottomrule
\end{tabular}
\end{center}

\vspace{0.5em}
\noindent\textbf{Días de la semana:}

\begin{center}
\begin{tabular}{lll}
\toprule
\textbf{Carácter} & \textbf{Pinyin} & \textbf{Día} \\
\midrule
星期一 & xīngqī yī & lunes \\
星期二 & xīngqī èr & martes \\
星期三 & xīngqī sān & miércoles \\
星期四 & xīngqī sì & jueves \\
星期五 & xīngqī wǔ & viernes \\
星期六 & xīngqī liù & sábado \\
星期日 / 星期天 & xīngqī rì / xīngqī tiān & domingo \\
\bottomrule
\end{tabular}
\end{center}

\newpage

\noindent\textbf{Meses del año:} El patrón es simple: número + 月 (yuè).

\begin{center}
\begin{tabular}{lll|lll}
\toprule
\textbf{Mes} & \textbf{Carácter} & \textbf{Pinyin} & \textbf{Mes} & \textbf{Carácter} & \textbf{Pinyin} \\
\midrule
Enero    & 一月  & yī yuè    & Julio      & 七月  & qī yuè \\
Febrero  & 二月  & èr yuè    & Agosto     & 八月  & bā yuè \\
Marzo    & 三月  & sān yuè   & Septiembre & 九月  & jiǔ yuè \\
Abril    & 四月  & sì yuè    & Octubre    & 十月  & shí yuè \\
Mayo     & 五月  & wǔ yuè    & Noviembre  & 十一月 & shíyī yuè \\
Junio    & 六月  & liù yuè   & Diciembre  & 十二月 & shí'èr yuè \\
\bottomrule
\end{tabular}
\end{center}

\subsection{1.2 Preguntar y hablar sobre acciones en el tiempo}

La estructura es: \textbf{tiempo + sujeto + verbo + complemento}. El tiempo va primero, antes del sujeto o inmediatamente después.

\begin{cuadroazul}
\textbf{Pregunta:} \quad 明天星期六，你去学校吗？\\
\textit{(Mañana es sábado, ¿vas a la escuela?)}\\[4pt]
\textbf{Respuesta afirmativa:} \quad 我去学校。\\[4pt]
\textbf{Pregunta de propósito:} \quad 你去学校做什么？\quad ¿Para qué vas a la escuela?\\[4pt]
\textbf{Respuesta:} \quad 我去学校看书。\quad Voy a la escuela a leer libros.
\end{cuadroazul}

\noindent\textbf{Vocabulario nuevo:}

\begin{center}
\begin{tabular}{lll}
\toprule
\textbf{Carácter} & \textbf{Pinyin} & \textbf{Significado} \\
\midrule
看书 & kàn shū & leer libros \\
学校 & xuéxiào & escuela / colegio \\
做什么 & zuò shénme & hacer qué / para qué \\
\bottomrule
\end{tabular}
\end{center}

\noindent\textbf{Nota gramatical clave:} En chino, dos verbos seguidos pueden indicar propósito (ir \textit{para} hacer algo). \mbox{去学校看书} = ir a la escuela [para] leer. No se necesita preposición extra.

%---------------------------------------------------------
\section{Contenido teórico — Lección 8: Deseos, comida y compras}
%---------------------------------------------------------

\subsection{2.1 Expresar deseo: 想 (xiǎng)}

\begin{cuadroazul}
\textbf{Estructura:} \quad 主语 + 想 + 动词 + 宾语\\
\textit{Sujeto + querer/desear + verbo + objeto}\\[6pt]
\textbf{Ejemplos:}\\
你想喝什么？\quad ¿Qué quieres beber?\\
我想喝茶。\quad Quiero beber té.\\
你想吃什么？\quad ¿Qué quieres comer?\\
我想吃米饭。\quad Quiero comer arroz.
\end{cuadroazul}

\newpage

\noindent\textbf{Vocabulario de comida y bebida:}

\begin{center}
\begin{tabular}{lll}
\toprule
\textbf{Carácter} & \textbf{Pinyin} & \textbf{Significado} \\
\midrule
喝 & hē & beber \\
吃 & chī & comer \\
茶 & chá & té \\
米饭 & mǐfàn & arroz (cocido) \\
中国菜 & zhōngguó cài & comida china \\
\bottomrule
\end{tabular}
\end{center}

\subsection{2.2 Actividades de la tarde, noche y mañana}

\begin{cuadroazul}
\textbf{Partes del día:}\\[4pt]
\begin{tabular}{lll}
上午 & shàngwǔ & mañana (antes del mediodía) \\
下午 & xiàwǔ & tarde \\
晚上 & wǎnshang & noche \\
今天上午 & jīntiān shàngwǔ & esta mañana \\
今天下午 & jīntiān xiàwǔ & esta tarde \\
今天晚上 & jīntiān wǎnshang & esta noche \\
明天上午 & míngtiān shàngwǔ & mañana por la mañana \\
\end{tabular}\\[6pt]
\textbf{Estructura:} tiempo + sujeto + 想 + verbo + objeto\\
下午你想做什么？\quad ¿Qué quieres hacer esta tarde?\\
下午我想去商店。\quad Esta tarde quiero ir a la tienda.
\end{cuadroazul}

\subsection{2.3 Preguntar precios: 多少钱 (duōshao qián)}

\begin{cuadroazul}
\textbf{Estructura:} \quad 这个/那个 + objeto + 多少钱？\\[4pt]
这个杯子多少钱？\quad ¿Cuánto cuesta este vaso?\\
二十八块。\quad Veintiocho yuanes.\\
那个被子多少钱？\quad ¿Cuánto cuesta ese edredón?\\
那个被子十八块钱。\quad Ese edredón cuesta dieciocho yuanes.\\[4pt]
\textbf{Nota:} 块 (kuài) es la forma coloquial de 元 (yuán). Ambas significan yuan.\\
这个 = este/esta \quad | \quad 那个 = ese/esa/aquel
\end{cuadroazul}

\noindent\textbf{Vocabulario de compras:}

\begin{center}
\begin{tabular}{lll}
\toprule
\textbf{Carácter} & \textbf{Pinyin} & \textbf{Significado} \\
\midrule
商店 & shāngdiàn & tienda \\
买 & mǎi & comprar \\
一个 & yī gè & uno (clasificador general) \\
被子 & bèizi & edredón / manta \\
杯子 & bēizi & vaso / taza \\
多少钱 & duōshao qián & ¿cuánto cuesta? \\
块 & kuài & yuan (forma coloquial) \\
\bottomrule
\end{tabular}
\end{center}

\subsection{2.4 Combinación con vocabulario anterior}

El profesor combina vocabulario ya conocido. Practica construir oraciones del tipo:

\begin{cuadroazul}
\textbf{Modelo:} \quad 星期六下午我想去朋友家吃中国菜。\\
\textit{El sábado por la tarde quiero ir a casa de mi amigo a comer comida china.}\\[6pt]
Elementos combinados: día + parte del día + 想 + destino + propósito
\end{cuadroazul}

%---------------------------------------------------------
\section{Actividades del día — Bloque 1 (aprox. 60 min)}
%---------------------------------------------------------
\textit{Escritura de caracteres. Trabaja en tu cuaderno.}

\subsection{Actividad 1 — Copia espaciada: días de la semana}

Para cada día de la semana: cópialo 3 veces mirando el modelo, cúbrelo, escríbelo 3 veces de memoria. Si cometes un error, repite el ciclo. No copies todos seguidos.

\begin{cuadrogris}
星期一 \quad 星期二 \quad 星期三 \quad 星期四 \quad 星期五 \quad 星期六 \quad 星期日
\end{cuadrogris}

Al terminar con los 7 días: escríbelos todos en orden sin mirar, una vez. Luego escríbelos en orden inverso (de domingo a lunes), una vez.

\subsection{Actividad 2 — Copia espaciada: vocabulario nuevo lecciones 7 y 8}

Aplica la misma técnica (3 mirando + 3 de memoria) a estos caracteres:

\begin{cuadrogris}
今天 \quad 昨天 \quad 明天 \quad 学校 \quad 看书 \quad 上午 \quad 下午 \quad 晚上\\[4pt]
想 \quad 喝 \quad 吃 \quad 茶 \quad 米饭 \quad 商店 \quad 买 \quad 杯子 \quad 被子
\end{cuadrogris}

\subsection{Actividad 3 — Escritura de fechas concretas}

Escribe en chino las siguientes fechas en tu cuaderno, incluyendo el día de la semana donde se indique. Usa la estructura completa.

\begin{enumerate}[leftmargin=1.5cm]
\item Hoy (domingo, 22 de junio)
\item Tu cumpleaños (mes + día + día de la semana si lo sabes)
\item La fecha de tu próxima clase de chino (martes 24 de junio)
\item El día de tu primer examen de chino (escríbelo como lo recuerdes)
\item Ayer y mañana con sus días de la semana correspondientes
\end{enumerate}

%---------------------------------------------------------
\section{Actividades del día — Bloque 2 (aprox. 60 min)}
%---------------------------------------------------------
\textit{Producción oral y escrita con los diálogos. Trabaja en tu cuaderno.}

\subsection{Actividad 4 — Dictado activo de los diálogos}

Cubre los diálogos de tu apunte. Escribe de memoria cada línea de los tres diálogos de la Lección 7, verificando después. Cada error: repite esa línea 5 veces.

Los tres diálogos a recuperar de memoria son:

\begin{cuadrogris}
\textbf{Diálogo 1:}\\
请问，今天几号？\\
今天九月一号。\\
今天星期几？\\
星期三。\\[6pt]
\textbf{Diálogo 2:}\\
昨天是几月几号？\\
昨天是八月三十一号，星期二。\\
明天呢？\\
明天是九月二号，星期四。\\[6pt]
\textbf{Diálogo 3:}\\
明天星期六，你去学校吗？\\
我去学校。\\
你去学校做什么？\\
我去学校看书。
\end{cuadrogris}

\subsection{Actividad 5 — Responde estas preguntas por escrito}

Responde en chino usando información real tuya. Escríbelas en tu cuaderno.

\begin{enumerate}[leftmargin=1.5cm]
\item 今天几号？
\item 今天星期几？
\item 昨天是几月几号？
\item 明天是几月几号，星期几？
\item 你的生日是几月几号？\quad (生日 = cumpleaños, shēngrì)
\item 明天你去学校吗？
\item 你去学校做什么？
\item 星期六你去学校吗？
\end{enumerate}

\subsection{Actividad 6 — Construcción libre de oraciones (Lección 7)}

Escribe en tu cuaderno 8 oraciones propias usando la estructura tiempo + sujeto + verbo + complemento. Usa solo vocabulario que ya conoces. Ejemplos del tipo que debes construir (no las copies, crea las tuyas):

\begin{cuadrogris}
星期二我去学校学习汉语。\\
明天下午我不去学校。\\
昨天是星期三，今天是星期四。
\end{cuadrogris}

%---------------------------------------------------------
\section{Actividades del día — Bloque 3 (aprox. 60 min)}
%---------------------------------------------------------
\textit{Lección 8: deseos, comida y compras.}

\subsection{Actividad 7 — Dictado activo de los diálogos de Lección 8}

Mismo procedimiento que la Actividad 4. Cubre y escribe de memoria:

\begin{cuadrogris}
\textbf{Diálogo 1:}\\
你想喝什么？\\
我想喝茶。\\
你想吃什么？\\
我想吃米饭。\\[6pt]
\textbf{Diálogo 2:}\\
下午你想做什么？\\
下午我想去商店。\\
你想买什么？\\
我想买一个被子。\\[6pt]
\textbf{Diálogo 3:}\\
你好！这个杯子多少钱？\\
二十八块。\\
那个杯子多少钱？\\
那个被子十八块钱。
\end{cuadrogris}

\subsection{Actividad 8 — Responde por escrito}

\begin{enumerate}[leftmargin=1.5cm]
\item 你想喝什么？
\item 你想吃什么？
\item 今天下午你想做什么？
\item 今天晚上你想做什么？
\item 明天上午你想做什么？
\item 你想买什么？
\item 这个杯子多少钱？（Inventa un precio）
\item 你想去商店买什么？
\end{enumerate}

\subsection{Actividad 9 — Construcción combinada: vocabulario antiguo + nuevo}

Escribe en tu cuaderno las siguientes oraciones \textbf{sin mirar referencia}. Usa el vocabulario de los exámenes 1 y 2 combinado con el nuevo:

\begin{enumerate}[leftmargin=1.5cm]
\item Di que el sábado por la tarde quieres ir a casa de tu amigo.
\item Di que quieres comer comida china con tu compañero de clase.
\item Pregunta a tu profesor qué quiere beber.
\item Di que tu mamá no habla chino pero quiere aprender.
\item Di que hoy es lunes y mañana quieres ir a la escuela a estudiar chino.
\item Di que ayer fuiste a la tienda y compraste un vaso.
\item Pregunta cuánto cuesta ese edredón.
\item Di que tu familia tiene seis personas y esta noche quieren comer arroz.
\end{enumerate}

%---------------------------------------------------------
\section{Actividades del día — Bloque 4 (aprox. 30--40 min)}
%---------------------------------------------------------
\textit{Revisión y repaso de exámenes anteriores.}

\subsection{Actividad 10 — Dictado de vocabulario del examen 1}

Pide a alguien que te dicte las palabras en español, o dítelas a ti mismo en voz alta con los ojos cerrados y escribe el carácter. Si estás solo, lee la palabra en español, cierra los ojos, escribe el carácter, y luego verifica.

\begin{cuadrogris}
Hola (formal) · Gracias · Perdón · De nada · Adiós · Nosotros · China (país) · Estadounidense · Nombre · Profesor
\end{cuadrogris}

\subsection{Actividad 11 — Dictado de vocabulario del examen 2}

\begin{cuadrogris}
Amigo · Hablar chino · Saber/poder · Casa · Este año · Compañero de clase · Mamá · Comida china · Seis miembros · 56 años
\end{cuadrogris}

\subsection{Actividad 12 — Reordena y traduce}

Sin mirar referencia, reordena las siguientes palabras para formar oraciones correctas y tradúcelas al español en tu cuaderno:

\begin{enumerate}[leftmargin=1.5cm]
\item 想 / 喝 / 我 / 茶 / 下午
\item 星期六 / 去 / 你 / 吗 / 商店
\item 做 / 明天 / 什么 / 你 / 想
\item 家 / 朋友 / 去 / 晚上 / 我 / 的 / 想
\item 中国菜 / 吃 / 同学 / 和 / 我 / 想 / 我的
\item 几号 / 昨天 / 是 / 几月
\end{enumerate}

\vspace{1em}
\hrule
\vspace{0.5em}

\begin{center}
\textcolor{gray}{\small \textbf{Al terminar hoy:} revisa los caracteres nuevos una última vez. Mañana el foco cambia a producción oral y repaso integrado.}
\end{center}

\newpage
%---------------------------------------------------------
\section*{\textcolor{rojosuave}{Respuestas — Actividades con respuesta única}}
%---------------------------------------------------------

\textbf{Actividad 3 — Fechas:}
\begin{enumerate}[leftmargin=1.5cm]
\item 今天是六月二十二号，星期日。
\item (Depende de tu fecha de nacimiento)
\item 六月二十四号，星期二。
\item (Depende de tu examen)
\item 昨天是六月二十一号，星期六。\quad 明天是六月二十三号，星期一。
\end{enumerate}

\textbf{Actividad 5 — Preguntas sobre el tiempo:}
\begin{enumerate}[leftmargin=1.5cm]
\item 今天六月二十二号。
\item 今天星期日。
\item 昨天是六月二十一号。
\item 明天是六月二十三号，星期一。
\item (Depende de tu cumpleaños)
\item 明天我不去学校。 / 明天我去学校。 (según tu caso)
\item 我去学校学习。/ 我去学校看书。
\item 星期六我不去学校。 / 星期六我去学校。
\end{enumerate}

\textbf{Actividad 8 — Respuestas con 想:}
\begin{enumerate}[leftmargin=1.5cm]
\item 我想喝茶。/ 我想喝水。(水 = agua, shuǐ)
\item 我想吃米饭。/ 我想吃中国菜。
\item 今天下午我想学习汉语。/ 今天下午我想去商店。
\item 今天晚上我想看书。
\item 明天上午我想去学校。
\item 我想买一个杯子。/ 我想买一个被子。
\item (Inventa un precio: 三十块 / 五十块, etc.)
\item 我想去商店买被子。/ 我想去商店买杯子。
\end{enumerate}

\textbf{Actividad 9 — Oraciones combinadas:}
\begin{enumerate}[leftmargin=1.5cm]
\item 星期六下午我想去朋友家。
\item 我想和同学吃中国菜。
\item 老师，你想喝什么？
\item 我妈妈不会说汉语，但是她想学。\quad (但是 = pero, dànshì — vocabulario de apoyo)
\item 今天是星期一，明天我想去学校学习汉语。
\item 昨天我去商店买了一个杯子。\quad (Nota: 了 marca pasado — lo verás próximamente)
\item 那个被子多少钱？
\item 我家有六口人，今天晚上我们想吃米饭。
\end{enumerate}

\textbf{Actividad 12 — Reordenamiento:}
\begin{enumerate}[leftmargin=1.5cm]
\item 下午我想喝茶。\quad \textit{Esta tarde quiero beber té.}
\item 星期六你去商店吗？\quad \textit{¿El sábado vas a la tienda?}
\item 明天你想做什么？\quad \textit{¿Qué quieres hacer mañana?}
\item 晚上我想去朋友的家。\quad \textit{Esta noche quiero ir a casa de mi amigo.}
\item 我想和我的同学吃中国菜。\quad \textit{Quiero comer comida china con mi compañero de clase.}
\item 昨天是几月几号？\quad \textit{¿Qué fecha fue ayer?}
\end{enumerate}

\end{document}